\uuid{IWjs}
\titre{Changements de variables}
\auteur{hanani ; mzari}
\organisation{exo7}
\licence{CC BY-NC-SA 4.0}
\datecreate{2020-01-20}
\niveau{L1}
\module{Analyse}
\chapitre{Calcul d'intégrales}
\sousChapitre{Autre}
\difficulte{3}
\theme{Calculs d'intégrales, changements de variables}
\exo7qcmid{lille-2.integrales.019}

\contenu{
\begin{question}
\qclassification{127.99}
\qtags{motcle=changements de variables}
On pose $\displaystyle I_1=\int _{0}^{\pi/2}\frac{\cos x\,\mathrm{d}x}{1+2\sin x}$, $\displaystyle I_2=\int _{0}^{\pi/2}\frac{\sin (2x)\,\mathrm{d}x}{1+2\sin x}$ et $I=I_1+I_2$. Parmi les affirmations suivantes, cocher celles qui sont vraies :
\begin{answers}
\good{$\displaystyle I=1$.}
\bad{$\displaystyle I_1=2\ln 3$.}
\good{$\displaystyle I_1=\frac{1}{2}\ln 3$.}
\bad{$\displaystyle I_2=1-2\ln 3$.}
\end{answers}
\vskip2mm
\begin{explanations}
Dans $I_2$, on écrit $\sin (2x)=2\sin x\cos x$. D'où $\displaystyle I_1+I_2=\int _{0}^{\pi/2}\cos x\,\mathrm{d}x=1$.
\vskip0mm
Avec $t=1+2\sin x\Rightarrow \mathrm{d}t=2\cos x\, \mathrm{d}x$, on obtient : $\displaystyle I_1=\int _{1}^{3}\frac{\mathrm{d}t}{2t}=\frac{1}{2}\ln 3$. Enfin, $I_2=1-I_1$.
\end{explanations}
\end{question}
}
